%% ===========================================================================
\section*{How to read this document}
\addcontentsline{toc}{section}{How to read this document}

\subsection*{Thesis}

Load-bearing, but not printed as a sentence in the paper:

% --- Previous thesis (commented out per Mohit: avoid \sout, comment old + add new) ---
% In inter-state crisis discourse, the dynamics that matter for harm and for
% moderation --- where hostility concentrates, which voices get suppressed,
% which topics crack an echo chamber open --- are properties of
% \emph{compositions} of annotation layers and community context, and are
% largely invisible to the single layer (toxicity) that platform moderation is
% actually built on.

% --- Version A (commented; superseded by A-2, which also covers RQ3) ---
% In inter-state crisis discourse, what a community contests --- and how the
% same narrative frame is received across communities --- is a property of the
% \emph{composition} of a narrative frame with its community context, not of any
% single comment scored on its own. Approaches that measure discourse one layer
% at a time, including the toxicity layer, largely miss this composition.

\begin{quote}
  In inter-state crisis discourse, what a community contests, how it receives the
  same narrative frame, and what happens when its members cross into an opposing
  community are all properties of the \emph{composition} of a narrative frame
  with its community context, not of any single comment scored on its own.
  Approaches that measure discourse one layer at a time, including the toxicity
  layer, largely miss this composition.
\end{quote}

\jun{Proposed thesis (Version A-2; Version A commented above). A-2 keeps A's two
changes from the team's original --- dropping ``harm''/``moderation'' (we run no
moderation system, and \texttt{controversiality} measures contested reception,
not harm) and re-centring on what we measured --- but widens the spine to cover
all three RQs instead of two. The three clauses map one-to-one:
``what a community contests'' = RQ1 (R8); ``how it receives the same frame'' =
RQ2 (R5); ``what happens when its members cross'' = RQ3 (R6). All three are
folded under the same unifying object (frame~$\times$~community composition,
invisible to single-comment scoring), which fixes the gap in Version A where RQ3
had no home in the spine. R7 (the spiral-of-silence null) stays out of the spine
by design and is carried as the C3 methodological result. Open for team
discussion before we align the abstract and \S1.4.}

\subsection*{Contributions}

Goes verbatim into \S1.4.

\begin{description}
\item[C1 (empirical).] A compositional analysis of 309,204 annotated comments
  across six India--Pakistan crises (2019--2025) and 14 subreddits, showing that
  (a)~crisis \emph{type} determines the frame--affect signature more than era or
  temporal proximity, (b)~communities render the same frame with systematically
  different affect and differentially \emph{reward} it, (c)~cross-community
  participation is antagonistic rather than bridging and is penalized regardless
  of rhetorical register, and (d)~contested discourse tracks frame--community
  mismatch rather than toxicity.

\item[C2 (framework and measurement).] An operationalization of compositional
  harm analysis for inter-state crisis discourse: a six-layer annotation scheme
  with an explicit gating structure, a set of composition operators (per-frame
  divergence, transition matrices, author-level compositional fingerprints) that
  recover effects no single layer exposes, and \textbf{a labeller-comparison
  protocol across commercial LLMs, open-source LLMs, fine-tuned encoders, and
  human coders} --- showing that no automated labeller reaches a human ceiling
  that is itself only $\kappa$ 0.45--0.68, and that the Perspective API is the
  weakest instrument in the comparison.

\item[C3 (methodological, negative).] A replicated demonstration that
  cluster-robust difference-in-differences over-rejects on crisis-discourse
  panels of this shape, and that the parallel-trends assumption fails
  structurally when crises cluster in time --- so causal designs on six-event,
  ten-subreddit corpora are \emph{uninformative}, not merely underpowered.
\end{description}